\section{The Four Approaches to Semantics Compared}
\label{sec:semantics-comparison}
\slidesrc{Blackboard: semantics B; Blackboard: semantics A}

We have now met four ways of giving meaning to words and sentences:
\textbf{distributional} semantics (Section~\ref{sec:distributional-semantics}),
\textbf{frame} semantics (Section~\ref{sec:frame-semantics}),
\textbf{model-theoretic} semantics (Section~\ref{sec:semantic-parsing}) and
\textbf{lexical} semantics (Section~\ref{sec:lexical-semantics}). In class the lecturer
summarised them in one table on the blackboard, comparing their main features, pros,
cons, applications and origins. The table is rewritten below exactly as it was written on
the board; the explanations follow.

\begin{nbblackboard}{The four kinds of semantics compared}{Blackboard: semantics B
(left half); Blackboard: semantics A (right half)}
\footnotesize
\renewcommand{\arraystretch}{1.25}
\begin{tabularx}{\linewidth}{@{}>{\bfseries\raggedright\arraybackslash}p{1.35cm}
  >{\raggedright\arraybackslash}X >{\raggedright\arraybackslash}X
  >{\raggedright\arraybackslash}X >{\raggedright\arraybackslash}X@{}}
\toprule
 & \textbf{Distributional} & \textbf{Frame} & \textbf{Model-theoretic} & \textbf{Lexical} \\
\midrule
Main\newline features &
  words represented as vectors; vectors' distance [measures how different the meanings
  are] &
  meaning of words: their role inside frames given by pre-compiled [frame resources] &
  semantics of words \& sentences given by their logical meaning, as terms, atoms,
  formulae, \ldots &
  listing relations among words composing the lexicon of a language \\
 &
  \textcolor{nbred}{$\bullet$ \textbf{statistical} $\bullet$} &
  \textcolor{nbred}{$\bullet$ \textbf{human-driven} $\bullet$} &
  \textcolor{nbred}{$\bullet$ \textbf{translation-based} $\bullet$} &
  \textcolor{nbred}{$\bullet$ \textbf{human-driven} $\bullet$} \\
\midrule
Pros &
  learned by software \newline {\color{nbgray}\itshape(a second line was erased)} &
  explainability; explicit &
  in very limited application domains it may work; mainly for \underline{statement}
  translation &
  lightweight creation; BabelNet, WordNet exist \\
\midrule
Cons &
  implicit, monolithic learning of the vector space (?); need of large corpus &
  human annotation; not portable &
  not ``very feasible'' in practice &
  human annotation \\
\midrule
Applic\-ations &
  disambiguate meaning (with some intermediate steps); comparison &
  disambiguate by checking whether a word fits one frame &
  querying databases; reasoning on the consistency of the sentence &
  information filtering, retrieval; helping DS \\
\midrule
Origins &
  Wittgenstein &
  Fillmore &
  Leibniz, Montague, \ldots &
  {\color{nbgray}\itshape(left empty)} \\
\bottomrule
\end{tabularx}

\smallskip
{\color{nbgray}\scriptsize Words in square brackets are cut off at the edge of the
photos and completed from context. ``DS'' = distributional semantics. The ``(?)'' is on
the board.}
\end{nbblackboard}

\begin{nbwarning}[how the two photos fit together]
The board was photographed in two halves. The first photo holds the row headings and the
\emph{Distributional} and \emph{Frame} columns; the second holds the
\emph{Model-theoretic} and \emph{Lexical} columns. The \emph{Frame} column runs across
the gap between two blackboard panels: ``human-dri\ldots'' + ``\ldots ven'', ``human
annotati\ldots'' + ``\ldots on'', and ``explainability'' is followed by ``explicit'' on
the next panel. So there are four columns, not five.
\end{nbwarning}

% ---------------------------------------------------------------------------
\subsection{Reading the table column by column}
\label{sec:semantics-comparison-columns}
\slidesrc{Blackboard: semantics B; Blackboard: semantics A}

\paragraph{Distributional semantics (statistical).}
\begin{itemize}[nosep]
  \item \emph{Main features}: every word is a vector of numbers learnt from the contexts
        in which it occurs; the distance (angle) between two vectors measures how
        different the meanings are (Section~\ref{sec:geometric-interpretation}). It is
        \textbf{statistical}: nobody writes the meanings by hand, they come out of counts
        over a corpus.
  \item \emph{Pro}: it is \textbf{learned by software}, automatically, for any word that
        appears often enough in the corpus.
  \item \emph{Cons}: the knowledge is \textbf{implicit}: a vector like
        $(0.031, 0.143, -0.243, \ldots)$ does not say \emph{why} two words are similar.
        The vector space is learnt \textbf{monolithically}, all at once from the whole
        corpus: one cannot easily correct a single word by hand (the ``(?)'' on the board
        signals this is a debatable point). And it needs a \textbf{large corpus}: rare
        words get unreliable vectors.
  \item \emph{Applications}: \textbf{disambiguating meaning}, but only ``with some
        intermediate steps'' (context vectors, clustering of senses:
        Section~\ref{sec:dsm-in-practice}); \textbf{comparison} of words (similarity,
        synonyms, nearest neighbours).
  \item \emph{Origins}: \textbf{Wittgenstein}, ``the meaning of a word is its use''
        (Section~\ref{sec:distributional-hypothesis}).
\end{itemize}

\paragraph{Frame semantics (human-driven).}
\begin{itemize}[nosep]
  \item \emph{Main features}: the meaning of a word is the \textbf{role} it plays inside
        a \textbf{frame} (a situation with participants, e.g.\ the commercial-transaction
        frame with buyer, seller, goods, money; Section~\ref{sec:frame-semantics}). The
        frames are \textbf{pre-compiled} by people (FrameNet, PropBank): it is
        \textbf{human-driven}.
  \item \emph{Pros}: \textbf{explainability}: a frame states explicitly that \emph{Cynthia}
        is the seller and \emph{the bike} the goods; the representation is
        \textbf{explicit}, readable by a human.
  \item \emph{Cons}: frames require \textbf{human annotation}, which is slow and costly;
        and they are \textbf{not portable}: frames written for one language or domain do
        not carry over to another.
  \item \emph{Applications}: \textbf{disambiguation by checking whether a word fits one
        frame}: \emph{bank} in ``I deposited money at the bank'' fits the frame of
        financial transactions, not the frame of rivers.
  \item \emph{Origins}: \textbf{Charles Fillmore} (1977), who introduced frames.
\end{itemize}

\paragraph{Model-theoretic semantics (translation-based).}
\begin{itemize}[nosep]
  \item \emph{Main features}: the meaning of words and sentences is their
        \textbf{logical meaning}: words become \textbf{terms} and predicates, simple
        sentences become \textbf{atoms}, complex sentences become \textbf{formulae},
        which can be evaluated as true or false in a model of the world
        (Section~\ref{sec:semantic-parsing}). It is \textbf{translation-based}: the
        sentence is \emph{translated} into logic.
  \item \emph{Pros}: in \textbf{very limited application domains} it may work (a database
        of US geography, a calendar, the blocks world), mainly for the translation of
        \textbf{statements} (declarative sentences); questions, commands, irony,
        metaphors are much harder.
  \item \emph{Cons}: \textbf{not ``very feasible'' in practice} for open, real language:
        translating arbitrary text into correct logic is extremely hard.
  \item \emph{Applications}: \textbf{querying databases} (natural-language interfaces:
        ``which states border Texas?'' becomes a query); \textbf{reasoning on the
        consistency} of a sentence (is it contradictory? does it follow from what we
        know?).
  \item \emph{Origins}: \textbf{Leibniz} (the dream of a universal logical language in
        which disputes are settled by calculation), \textbf{Montague} (natural language
        can be treated like a formal language, compositionally), \ldots
\end{itemize}

\paragraph{Lexical semantics (human-driven).}
\begin{itemize}[nosep]
  \item \emph{Main features}: it \textbf{lists the relations} among the words that make
        up the lexicon of a language: synonymy, antonymy, hyponymy, meronymy, \ldots
        (Section~\ref{sec:lexical-relations}). It is \textbf{human-driven}: lexicographers
        write the relations.
  \item \emph{Pros}: \textbf{lightweight creation} (much simpler than frames or logic:
        a relation is just a pair of words with a label), and large resources already
        \textbf{exist}: \textbf{WordNet} (Section~\ref{sec:wordnet}) and
        \textbf{BabelNet}, its multilingual extension.
  \item \emph{Cons}: \textbf{human annotation}, like frames.
  \item \emph{Applications}: \textbf{information filtering} and \textbf{retrieval} (e.g.\
        expanding a query with synonyms and hyponyms); \textbf{helping distributional
        semantics} (``helping DS''): e.g.\ WordNet senses to label the clusters of a
        DSM, or lexical relations to evaluate and correct vectors.
  \item \emph{Origins}: left empty on the board.
\end{itemize}

\begin{nbintuition}[the big picture]
Read the table along the ``$\bullet$'' row. One approach is \textbf{statistical}
(meaning is \emph{learnt} from data), two are \textbf{human-driven} (meaning is
\emph{written} by people), one is \textbf{translation-based} (meaning is a
\emph{logical formula}). The pros and cons follow from this one choice:
\begin{itemize}[nosep]
  \item statistical $\Rightarrow$ cheap, automatic, covers everything, but implicit and
        hungry for data;
  \item human-driven $\Rightarrow$ explicit and explainable, but costly to build and hard
        to port;
  \item translation-based $\Rightarrow$ precise, supports reasoning, but works only in
        small, controlled domains.
\end{itemize}
This is the same trade-off as between the eras of NLP (Section~\ref{sec:nlp-history}):
rule-based (human-driven) vs.\ statistical and neural (learned) methods.
\end{nbintuition}

\begin{nbsummary}[the four semantics]
\begin{itemize}[nosep]
  \item \textbf{Distributional}: vectors, distance; statistical; learned automatically;
        implicit, needs a large corpus; disambiguation (indirectly), comparison;
        Wittgenstein.
  \item \textbf{Frame}: role of words inside pre-compiled frames; human-driven;
        explainable, explicit; human annotation, not portable; disambiguation by frame
        fitting; Fillmore.
  \item \textbf{Model-theoretic}: logical meaning (terms, atoms, formulae);
        translation-based; works in limited domains, for statements; not feasible in
        general; database querying, consistency reasoning; Leibniz, Montague.
  \item \textbf{Lexical}: relations among the words of the lexicon; human-driven;
        lightweight, WordNet and BabelNet exist; human annotation; information
        filtering/retrieval, helping DS.
\end{itemize}
\end{nbsummary}
